\documentclass[11pt,a4paper]{article}

\usepackage{iftex}
\ifPDFTeX
  \usepackage[utf8]{inputenc}
  \usepackage[T5,T1]{fontenc}
  \usepackage[english,vietnamese]{babel}
\else
  \usepackage{fontspec}
  \usepackage{polyglossia}
  \setdefaultlanguage{vietnamese}
  \setotherlanguage{english}
\fi
\usepackage[a4paper,margin=2.15cm,headheight=15pt]{geometry}
\usepackage{microtype}
\usepackage{graphicx}
\usepackage{array,booktabs,tabularx}
\usepackage{amsmath}
\usepackage{enumitem}
\usepackage{xcolor}
\usepackage{tikz}
\usetikzlibrary{arrows.meta,positioning,fit,calc}
\usepackage{hyperref}
\usepackage{xurl}
\usepackage{fancyhdr}
\usepackage{titlesec}

\definecolor{ink}{HTML}{19324D}
\definecolor{teal}{HTML}{087E8B}
\definecolor{sky}{HTML}{EAF4F7}
\definecolor{mint}{HTML}{EAF5F0}
\definecolor{amber}{HTML}{FFF4D6}
\definecolor{linegray}{HTML}{D7E0E7}
\definecolor{muted}{HTML}{526273}
\hypersetup{
  colorlinks=true,
  linkcolor=teal,
  citecolor=teal,
  urlcolor=teal,
  pdftitle={Feasibility kinh doanh và kiến trúc AI ghép tutor--student tự kỷ},
  pdfauthor={SPM Mobile},
  pdfsubject={Nghiên cứu sản phẩm, kinh doanh và AI hỗ trợ điều phối giáo dục}
}
\urlstyle{same}
\setlength{\parindent}{0pt}
\setlength{\parskip}{5pt}
\setlist[itemize]{leftmargin=1.35em,itemsep=2pt,topsep=3pt}
\setlist[enumerate]{leftmargin=1.55em,itemsep=3pt,topsep=3pt}
\renewcommand{\arraystretch}{1.22}
\newcolumntype{Y}{>{\raggedright\arraybackslash}X}
\newcommand{\vn}[1]{\textcolor{ink}{#1}}
\newcommand{\assumption}{\textcolor{teal}{\textbf{GIẢ ĐỊNH MINH HỌA}}}
\newcommand{\notmeasured}{\textcolor{muted}{\textit{Chưa đo trong pilot}}}

\pagestyle{fancy}
\fancyhf{}
\lhead{\small\textcolor{muted}{SPM Mobile \textbar\ Nghiên cứu sản phẩm}}
\rhead{\small\textcolor{muted}{Bản nghiên cứu mở rộng}}
\cfoot{\small\textcolor{muted}{\thepage}}
\titleformat{\section}{\Large\bfseries\color{ink}}{\thesection}{0.65em}{}
\titleformat{\subsection}{\large\bfseries\color{teal}}{\thesubsection}{0.65em}{}
\titleformat{\subsubsection}{\normalsize\bfseries\color{ink}}{\thesubsubsection}{0.65em}{}

\begin{document}

\begin{center}
{\large\textcolor{teal}{NGHIÊN CỨU FEASIBILITY SẢN PHẨM VÀ KINH DOANH}}\\[7pt]
{\LARGE\bfseries\color{ink} Kết nối học sinh tự kỷ\\ với gia sư phù hợp}\\[7pt]
{\large Mô hình ghép nối AI-native, coordinator kiểm soát shortlist}\\[12pt]
\rule{0.82\textwidth}{0.7pt}\\[7pt]
\textcolor{muted}{SPM Mobile \quad|\quad Việt Nam \quad|\quad Chốt bằng chứng: 02/10/2026}
\end{center}

\vspace{5pt}
\noindent
\colorbox{sky}{\parbox{\dimexpr\textwidth-2\fboxsep\relax}{
\textbf{\textcolor{ink}{Kết luận tạm thời.}}
Có cơ sở để thử một dịch vụ ghép nối quy mô nhỏ, do coordinator điều phối và
gia đình/học sinh xác nhận. Bằng chứng hiện tại chưa đủ để kết luận quy mô thị
trường trả phí hoặc tính khả thi khi mở rộng. Vì vậy, nên kiểm chứng nhu cầu,
nguồn tutor và unit economics bằng pilot thủ công trước; AI chỉ nên đưa ra danh
sách gợi ý có giải thích để coordinator duyệt, không tự quyết định ai được ghép
với ai.
}}

\newpage
\tableofcontents
\newpage

\section{Mục tiêu và kết luận điều hành}

Nghiên cứu này đánh giá giả thuyết: một ứng dụng có thể giúp gia đình tìm gia
sư theo đúng môn học, lịch, cách học và hỗ trợ mà từng học sinh cần; coordinator
đảm bảo an toàn và duyệt đề xuất; AI giảm thời gian tìm kiếm bằng cách xếp hạng
các hồ sơ đã qua điều kiện cơ bản. Đây là dịch vụ giáo dục và kết nối, không
phải công cụ chẩn đoán tự kỷ hay thay thế chuyên gia lâm sàng.

\begin{tabularx}{\textwidth}{>{\bfseries\color{ink}}p{3.25cm}Y}
\toprule
Khía cạnh & Đánh giá tại thời điểm lập báo cáo \\
\midrule
Nhu cầu xã hội & Có tín hiệu về rào cản giáo dục hòa nhập cho trẻ khuyết tật nói
chung tại Việt Nam. Các số liệu này không riêng cho tự kỷ và không chứng minh
nhu cầu mua gia sư. \\
Dữ liệu quy mô thị trường & Chưa có ước lượng đủ chắc cho số học sinh tự kỷ trong
độ tuổi học môn phổ thông, số gia đình tìm gia sư, khả năng chi trả hoặc số tutor
đủ năng lực. Không suy TAM từ tỷ lệ hiện mắc trẻ nhỏ. \\
Khả thi sản phẩm & Có thể kiểm chứng qua mô hình concierge: ghép thủ công, quy
mô hẹp, theo dõi mức chấp nhận, độ phù hợp, an toàn và thời gian coordinator. \\
Khả thi thương mại & Chưa xác nhận. Mô hình hoa hồng có thể chịu áp lực chi phí
điều phối; hợp đồng với trung tâm/trường hoặc tài trợ tiếp cận cần được thử song
song. \\
AI & Chỉ nên bật sau khi có hồ sơ tutor đã xác minh, consent, dữ liệu phù hợp và
quy trình giám sát. Human coordinator giữ quyền duyệt, từ chối và ghép lại. \\
\bottomrule
\end{tabularx}

\subsection{Quyết định đề xuất}

\textbf{Tiếp tục sang giai đoạn discovery và pilot có kiểm soát.} Chưa đầu tư
vào mô hình AI tự động hoặc tuyên bố đây là marketplace đã chứng minh được khả
năng mở rộng. Thứ tự nên là:
\begin{enumerate}
  \item Phỏng vấn gia đình, tutor và tổ chức giáo dục để xác nhận vấn đề, cách
  họ tìm nhau hiện nay, tiêu chí phù hợp, niềm tin và người trả tiền.
  \item Thử dịch vụ ghép thủ công với coordinator; ghi lại lý do được nhận,
  từ chối, đổi tutor và chi phí phục vụ từng trường hợp.
  \item Chỉ sau khi quy trình và dữ liệu đủ tốt, chạy AI ở chế độ shadow
  (xếp hạng nội bộ, chưa tác động lựa chọn), rồi đánh giá an toàn và độ phù hợp.
\end{enumerate}

\section{Nhu cầu và bằng chứng thị trường}

\subsection{Số liệu đã kiểm chứng}

\par\medskip\noindent
\textbf{Tự kỷ: bằng chứng trực tiếp tại Việt Nam.}\par\smallskip\noindent
\begin{tabularx}{\textwidth}{p{3.15cm}p{5.25cm}Y}
\toprule
\textbf{Nguồn và phạm vi} & \textbf{Số liệu được báo cáo} &
\textbf{Cách dùng và giới hạn} \\
\midrule
Khảo sát quần thể 7 tỉnh/thành, 2017--2018 \cite{vui2022} &
40.243 trẻ 18--30 tháng tham gia (94,5\% số được mời); 305 ca xác nhận sau
quy trình sàng lọc M-CHAT và đánh giá DSM-IV; tỷ lệ 0,758\%, xấp xỉ 1/132. &
Khảo sát đại diện quốc gia cho nhóm tuổi 18--30 tháng theo bài báo; không đo
học sinh tuổi đi học, nhu cầu gia sư hay khả năng chi trả. \\
\addlinespace
Nghiên cứu trường mầm non miền Nam, công bố 2024 \cite{tran2024} &
9.397 trẻ 24--72 tháng ở 41 trường; bác sĩ tâm thần nhi đánh giá theo DSM-5;
tỷ lệ 1,2\%, xấp xỉ 1/84. &
Một nghiên cứu khu vực, không phải ước lượng toàn quốc. Khác tuổi, địa bàn và
quy trình với khảo sát 2017--2018; không gộp hai tỷ lệ. \\
\addlinespace
WHO, ước tính toàn cầu \cite{who2025} &
Ước tính khoảng 1/127 người trên thế giới có tự kỷ vào năm 2021. &
Chỉ để đặt bối cảnh toàn cầu. WHO lưu ý khả năng và nhu cầu khác nhau, có thể
thay đổi; không áp tỷ lệ này vào dân số Việt Nam để suy ra số khách hàng. \\
\addlinespace
CDC/ADDM, Hoa Kỳ, năm giám sát 2022 \cite{cdc2025} &
32,2/1.000 trẻ 8 tuổi (khoảng 1/31) trên 16 địa điểm; địa điểm dao động
9,7--53,1/1.000. &
Định nghĩa giám sát kết hợp hồ sơ chẩn đoán, giáo dục đặc biệt và mã ICD.
Khác hệ thống phát hiện và bối cảnh Việt Nam; không chuyển thành tỷ lệ Việt Nam. \\
\bottomrule
\end{tabularx}

\vspace{5pt}
\textbf{Giáo dục, khả năng tiếp cận và mức sẵn sàng số: chủ yếu là số liệu
khuyết tật nói chung, không riêng tự kỷ.}\par\smallskip\noindent
\begin{tabularx}{\textwidth}{p{3.15cm}p{5.25cm}Y}
\toprule
\textbf{Nguồn và phạm vi} & \textbf{Số liệu được báo cáo} &
\textbf{Cách dùng và giới hạn} \\
\midrule
Điều tra người khuyết tật Việt Nam 2023 (VDS 2023) \cite{gso2023} &
Tỷ lệ khuyết tật: 1,98\% ở trẻ 2--17 tuổi; 1,51\% ở trẻ 5--17 tuổi; 1,56\%
ở trẻ 5--15 tuổi. &
Đây là các phân nhóm tuổi và khái niệm khuyết tật của khảo sát, không phải
tỷ lệ tự kỷ. Không nhân dân số học sinh với các tỷ lệ này để tính TAM tự kỷ. \\
\addlinespace
VDS 2023, đi học đúng tuổi \cite{gso2023} &
Tiểu học: trẻ khuyết tật 68,1\%, không khuyết tật 95,2\% (chênh 27,1 điểm
phần trăm); THCS: 53,0\% và 90,0\% (chênh 37 điểm); THPT: trẻ khuyết tật
30,8\%, thấp hơn nhóm không khuyết tật 45,8 điểm. &
Chỉ số là \emph{đi học đúng tuổi}, không phải tỷ lệ nhập học chung. Gợi ý cần
đo sự phù hợp và khả năng tiếp cận; không chứng minh gia đình cần hoặc mua tutor. \\
\addlinespace
VDS 2023, truy cập công nghệ thông tin \cite{gso2023} &
Người khuyết tật từ 6 tuổi có truy cập Internet trong 3 tháng: 33,6\%, so
với 83,7\% ở người không khuyết tật; sở hữu điện thoại di động lần lượt
53,7\% và 89,2\%. &
Khoảng cách này yêu cầu app nhẹ, dùng được với băng thông thấp và có kênh
hỗ trợ ngoài app. Không coi tỷ lệ của toàn nhóm khuyết tật là riêng trẻ tự kỷ. \\
\addlinespace
VDS 2023, trường có học sinh khuyết tật \cite{gso2023} &
79,6\% trường có học sinh khuyết tật, so với 71,4\% năm 2016; mức vùng
được báo cáo từ 58,8\% tại Đồng bằng sông Cửu Long đến 90,2\% tại Trung du
và miền núi phía Bắc. &
Cho thấy có mặt học sinh khuyết tật trong trường học nhưng không cho biết
mật độ tự kỷ, nguồn tutor hay nhu cầu trả phí tại từng địa phương. \\
\bottomrule
\end{tabularx}

\vspace{5pt}
\textbf{Hạ tầng số và chương trình hỗ trợ: không đồng nghĩa với khả năng tiếp
cận app của từng gia đình.}\par\smallskip\noindent
\begin{tabularx}{\textwidth}{p{3.15cm}p{5.25cm}Y}
\toprule
\textbf{Nguồn và phạm vi} & \textbf{Số liệu được báo cáo} &
\textbf{Cách dùng và giới hạn} \\
\midrule
UNICEF, kết quả hỗ trợ hòa nhập tại Việt Nam 2024 \cite{unicef2025inclusive} &
Báo cáo nêu hơn 400.000 trẻ khuyết tật chưa đến trường; 3\% trường tiếp cận
được về mặt thể chất; 1/10 trường THCS có giáo viên được đào tạo về giáo dục
hòa nhập. &
Số liệu UNICEF về trẻ khuyết tật nói chung; có thể dùng để mô tả bối cảnh hệ
thống, không phải số trẻ tự kỷ tiềm năng cho sản phẩm. \\
\addlinespace
UNESCO, AI Readiness Assessment Việt Nam 2025 \cite{unescoai2025} &
Khoảng 90\% trường kết nối Internet cho mục đích giảng dạy; tỷ lệ máy tính
xấp xỉ 1 máy/3 học sinh; gần 40\% trường thiếu học liệu số. &
Đây là chỉ số trường học, không phải kết nối hộ gia đình. Kết nối Internet
không đồng nghĩa có thiết bị riêng, nội dung tiếp cận được hoặc giáo viên sẵn
sàng dùng app. \\
\addlinespace
UNICEF--MOET, chương trình chuyển đổi số 2022--2026 \cite{unescoedtech2026} &
UNESCO SDG4 Hub ghi nhận 260.000 giáo viên/cán bộ quản lý được tập huấn,
4 triệu trẻ hưởng lợi trực tiếp (6\% là trẻ khuyết tật), triển khai cùng
10 Sở GD\&ĐT; ngân sách được nêu là 2,3 triệu USD. &
Đây là quy mô và kết quả chương trình, không phải số người dùng app này hay
bằng chứng AI matching đã nâng kết quả học tập. Phải giữ nhãn “báo cáo chương
trình”, không diễn giải thành thị trường. \\
\addlinespace
World Bank, dự án hỗ trợ giáo dục trẻ điếc \cite{worldbankdeaf2023} &
1.929 học sinh điếc ở 20/68 tỉnh, thành được dự án hỗ trợ; 96\% đạt bài kiểm
tra cuối kỳ so với mục tiêu dự án 50\%; có 150 video bài học và 4.000 mục từ
điển ngôn ngữ ký hiệu. &
Kết quả của gói tổng hợp gồm đào tạo, phương pháp, học liệu và hỗ trợ sư phạm;
không thể quy kết riêng cho công nghệ, và không suy rộng sang tự kỷ. Gợi ý
thiết kế phải kết hợp tutor, nội dung và điều phối. \\
\bottomrule
\end{tabularx}

\subsection{Diễn giải đúng cho quyết định kinh doanh}

Nguồn Việt Nam về tự kỷ hiện có trong desk research là nghiên cứu ở trẻ nhỏ;
không có số liệu quốc gia mới, được xác nhận ở đây cho số học sinh tự kỷ tuổi
đi học đang cần gia sư. Số liệu khuyết tật của GSO, số liệu UNICEF về trường
học, và tỷ lệ toàn cầu của WHO có đơn vị phân tích khác nhau. Chúng cho biết
bối cảnh và các rào cản cần thiết kế quanh sản phẩm, không phải số khách hàng.
Tỷ lệ hiện mắc cũng không đồng nghĩa với nhu cầu học thêm, sự sẵn sàng chia sẻ
dữ liệu, khả năng chi trả hay ý định mua.

\textbf{Kết luận thị trường:} có lý do để kiểm chứng vấn đề và thiết kế dịch vụ
bao trùm, nhưng hiện chưa có căn cứ để công bố con số TAM/SAM/SOM đáng tin cậy.
Trước khi quyết định mở rộng, cần dữ liệu địa bàn mục tiêu từ gia đình, tutor
đã xác minh và đối tác có thể giới thiệu hai phía.

\subsection{Cách tính quy mô thị trường mà không thổi phồng}

\begin{align*}
N^* &= \min(N_F,N_T,N_C),\\
R_{\text{annual}} &= N^* \times M_{\text{paid}} \times R_{\text{net/month}}.
\end{align*}

Trong đó $N_F$ là số gia đình có nhu cầu và đồng ý; $N_T$ là suất tutor đạt
chuẩn; $N_C$ là suất coordinator có thể theo dõi; $M_{\text{paid}}$ là số
tháng trả phí bình quân; $R_{\text{net/month}}$ là doanh thu ròng mỗi tháng.
Giá trị phục vụ được $N^*$ bị giới hạn bởi phía nhỏ nhất trong ba năng lực.
Các biến cần đo tại địa bàn: số yêu cầu thật theo môn/lớp/lịch/hình thức; tỷ
lệ qua sàng lọc điều kiện dịch vụ; tỷ lệ gia đình đồng ý thử; tỷ lệ tutor đã
xác minh nhận ca; tỷ lệ hoàn tất buổi đầu; thời lượng duy trì; và mức phí ròng
sau hoàn tiền, hủy buổi, thanh toán và điều phối. Không lấy 0,758\% ở trẻ
18--30 tháng hoặc tỷ lệ khuyết tật nói chung làm đầu vào cho công thức này.

\begin{tabularx}{\textwidth}{p{3.0cm}p{6.3cm}Y}
\toprule
\textbf{Tầng ước lượng} & \textbf{Định nghĩa dùng trong nghiên cứu} &
\textbf{Nguồn dữ liệu phải thu} \\
\midrule
TAM lý thuyết & Gia đình/học sinh thuộc địa bàn và độ tuổi dịch vụ, có mục
tiêu học tập phù hợp với dịch vụ ghép tutor. & Số học sinh theo địa bàn/tuổi,
khảo sát nhu cầu; dữ liệu tự nguyện từ đối tác; không suy từ prevalence. \\
SAM khả dụng & Phần TAM có thể tiếp cận qua đối tác/tutor, lịch, môn, phương
thức và yêu cầu hỗ trợ hiện có. & Danh sách tutor xác minh, lịch trống, khả
năng phục vụ tại chỗ/online, địa bàn và yêu cầu mà tutor nhận. \\
SOM giai đoạn pilot & Số ca app và coordinator có thể tiếp nhận an toàn trong
kỳ pilot và có khả năng được tutor nhận. & Yêu cầu thực nhận, phút điều phối,
số tutor nhận việc và kết quả buổi đầu; đo theo tuần. \\
\bottomrule
\end{tabularx}

\section{Khách hàng, giá trị và mô hình doanh thu}

\subsection{Các bên trong dịch vụ}

\begin{itemize}
  \item \textbf{Người học:} học sinh tự kỷ hoặc học sinh có nhu cầu hỗ trợ học
  tập tương tự; không bắt buộc cung cấp chẩn đoán nếu không cần cho việc học và
  gia đình không muốn chia sẻ.
  \item \textbf{Người đồng hành/quyết định:} phụ huynh hoặc người giám hộ; cùng
  học sinh phù hợp độ tuổi và khả năng tham gia xác định mục tiêu, cách giao
  tiếp và lựa chọn gia sư.
  \item \textbf{Tutor:} gia sư theo môn/khối lớp, công khai kinh nghiệm và phạm
  vi chuyên môn; hồ sơ phải được kiểm tra trước khi nhận đề xuất.
  \item \textbf{Coordinator:} nhân sự chịu trách nhiệm tiếp nhận nhu cầu, xem
  xét hồ sơ, kiểm tra phù hợp/an toàn, giải thích đề xuất, xin xác nhận, theo
  dõi buổi đầu và hỗ trợ đổi tutor.
  \item \textbf{Đối tác:} trung tâm giáo dục hòa nhập, trường học, tổ chức hỗ
  trợ gia đình hoặc mạng lưới tutor; có thể vừa giới thiệu cung-cầu vừa giúp
  bảo đảm tiêu chuẩn vận hành.
\end{itemize}

\subsection{Giá trị cần kiểm chứng}

\begin{tabularx}{\textwidth}{p{3.0cm}Y Y}
\toprule
\textbf{Bên liên quan} & \textbf{Giá trị giả thuyết} & \textbf{Bằng chứng cần
thu trong pilot} \\
\midrule
Gia đình/học sinh & Ít phải tự dò tìm; được xem lý do tutor phù hợp; có thể
nói rõ cách học, mục tiêu và điều cần tránh. &
Thời gian từ đăng ký đến đề xuất; tỷ lệ chấp nhận; cảm nhận an toàn; số lần
phải đổi tutor. \\
Tutor & Nhận yêu cầu có bối cảnh học tập rõ ràng, lịch và mục tiêu phù hợp
trước khi nhận lớp. &
Tỷ lệ tutor nhận đề xuất; tỷ lệ hoàn thành buổi đầu; lý do từ chối; thời gian
chuẩn bị. \\
Trung tâm/trường & Coordinator gom nhu cầu, hồ sơ tutor và theo dõi ghép nối
thành một quy trình có thể xem lại. &
Số giờ điều phối tiết kiệm/tăng; số gia đình được phục vụ; khả năng ký hợp
đồng hoặc tài trợ. \\
\bottomrule
\end{tabularx}

\subsection{Mô hình doanh thu nên thử}

\begin{enumerate}
  \item \textbf{B2B2C với trung tâm hoặc đối tác:} phí dịch vụ/tháng hoặc phí
  cho mỗi ca đang được điều phối. Đây là giả thuyết phù hợp để thử đầu tiên
  vì có thể tận dụng mạng lưới tutor, chuyên môn và niềm tin sẵn có.
  \item \textbf{Hoa hồng buổi học:} nền tảng thu một phần giao dịch sau khi
  tutor và gia đình đã xác nhận. Cần kiểm tra khả năng chi trả, phí thanh toán,
  hủy buổi, hoàn tiền và chi phí coordinator.
  \item \textbf{Hỗ trợ tiếp cận:} tài trợ hoặc suất học bổng từ tổ chức/trung
  tâm có thể giúp giảm rào cản giá; không nên dựa vào mô hình này nếu chưa có
  đối tác cam kết.
\end{enumerate}

\textbf{Khuyến nghị:} không lấy quảng cáo hoặc bán dữ liệu làm nguồn doanh thu.
Trong giai đoạn đầu, thử hợp đồng với một đối tác và hoa hồng minh bạch; không
thu phí phụ huynh cho một thuật toán chưa được đánh giá.

\subsection{Phân khúc và kênh kiểm chứng}

\begin{tabularx}{\textwidth}{p{2.8cm}p{5.5cm}Y}
\toprule
\textbf{Phân khúc} & \textbf{Nhu cầu giả thuyết} &
\textbf{Cần hỏi/đo trước khi kết luận} \\
\midrule
Gia đình và học sinh & Tìm tutor đúng môn, lịch, cách giao tiếp và điều chỉnh
đã chọn; muốn được biết vì sao một người được đề xuất. &
Cách tìm hiện tại; thời gian, chi phí và lần đổi tutor; ai quyết định và ai
trả tiền; thông tin nào không muốn chia sẻ; học sinh muốn tham gia thế nào. \\
Tutor & Nhận ca có yêu cầu dạy và lịch rõ ràng; biết hỗ trợ nào được kỳ vọng;
có điểm liên hệ khi cần đổi lịch/trao đổi. &
Môn/lớp có thể dạy; lịch và địa bàn nhận; kinh nghiệm có thể kiểm chứng; cách
được giới thiệu; điều kiện từ chối; phút chuẩn bị và hỗ trợ cần từ coordinator. \\
Trung tâm/trường/tổ chức & Điều phối nguồn tutor, xác minh quy trình và giúp
gia đình tiếp cận dịch vụ. &
Ai có thẩm quyền giới thiệu; nghĩa vụ bảo vệ trẻ; cách thanh toán/hợp đồng;
nguồn tutor thật; yêu cầu lưu trữ và báo cáo; có ngân sách hay cơ chế tài trợ. \\
Coordinator & Giải quyết trường hợp thiếu thông tin, ngoại lệ, an toàn và
đổi ghép; theo dõi lý do hệ thống đề xuất. &
Số phút/case; checklist xét duyệt; tải ca theo tuần; tỷ lệ override; loại lý
do thiếu tutor; thời gian phản hồi và quy trình sự cố. \\
\bottomrule
\end{tabularx}

\textit{Đây là giả thuyết phân khúc, không phải phân loại đã được phỏng vấn
hoặc chứng minh có sẵn lòng chi trả. Không gắn mọi học sinh tự kỷ với cùng một
nhu cầu, kiểu học hay môi trường học.}

\subsection{Kịch bản hòa vốn minh họa}

Bảng dưới đây không phải giá thị trường Việt Nam. Đây là phép tính nhạy cảm để
biết dữ liệu nào phải hỏi trong phỏng vấn. Tất cả đầu vào ngoài công thức đều
là giả định, chưa phải báo giá hay chi phí thực tế.

\begin{center}
\small
\begin{tabular}{p{2.25cm}rrrrr}
\toprule
\textbf{Giá/buổi} & \textbf{Hoa hồng} &
\shortstack{\textbf{Doanh thu/tháng}\\\textbf{mỗi học sinh}} &
\shortstack{\textbf{Đóng góp sau}\\\textbf{biến phí}} &
\shortstack{\textbf{Số học sinh}\\\textbf{hòa vốn}} &
\textbf{Kịch bản}\\
\midrule
200.000 đ & 10\% & 80.000 đ & 50.000 đ & 400 & Thấp \\
200.000 đ & 15\% & 120.000 đ & 90.000 đ & 223 & Thấp \\
200.000 đ & 20\% & 160.000 đ & 130.000 đ & 154 & Thấp \\
\addlinespace
350.000 đ & 10\% & 140.000 đ & 110.000 đ & 182 & Cơ sở \\
350.000 đ & 15\% & 210.000 đ & 180.000 đ & 112 & Cơ sở \\
350.000 đ & 20\% & 280.000 đ & 250.000 đ & 80 & Cơ sở \\
\addlinespace
500.000 đ & 10\% & 200.000 đ & 170.000 đ & 118 & Cao \\
500.000 đ & 15\% & 300.000 đ & 270.000 đ & 75 & Cao \\
500.000 đ & 20\% & 400.000 đ & 370.000 đ & 55 & Cao \\
\bottomrule
\end{tabular}
\normalsize
\end{center}

\small
\textit{Tất cả là giả định minh họa, không phải giá hay chi phí thị trường:
4 buổi/tháng/học sinh; biến phí 30.000 đ/học sinh/tháng; chi phí cố định
20 triệu đ/tháng. Doanh thu = giá/buổi $\times$ 4 $\times$ hoa hồng;
đóng góp = doanh thu trừ biến phí; hòa vốn = làm tròn lên
20 triệu/đóng góp. Bảng nhạy cảm đồng thời theo 3 giá buổi (200.000,
350.000, 500.000 đ) và 3 tỷ lệ (10\%, 15\%, 20\%). Chưa tính thuế, hoàn
tiền, chi phí thu hút khách hàng, số ca không hoạt động hoặc chi phí nhân sự
tăng thêm khi tải coordinator vượt công suất. Vì vậy, bảng không chứng minh
được khả năng sinh lời.}
\normalsize

\textbf{Công thức năng lực điều phối:}
\[
  \text{số ca theo dõi tối đa/tháng}
  = \frac{\text{giờ coordinator khả dụng/tháng}}
  {\text{giờ onboarding + matching + follow-up trung bình mỗi ca}}.
\]
Hai đầu vào này chưa được đo. Nếu coordinator phải tăng giờ hoặc tuyển thêm
người khi số ca tăng, chi phí không còn cố định; điểm hòa vốn ở bảng trên sẽ
tăng. Pilot phải ghi phút thực tế cho từng hoạt động thay vì dùng một chi phí
điều phối bình quân giả định.

\section{Sản phẩm và mức độ sẵn sàng của app hiện tại}

Tài liệu kiến trúc của repository mô tả Flutter mobile có đăng nhập và điều
hướng theo role, gồm student, lecturer và coordinator; phần matching tutor
hiện chỉ là utility luật đơn giản trong React web, chưa thấy được gọi từ luồng
chính. Source cũng chưa có AI matching endpoint/worker hoặc module lưu kết quả
matching \cite{spmarchitecture2026}. Do đó:

\begin{itemize}
  \item \textbf{Tận dụng được:} Flutter app, xác thực và khung điều hướng theo
  role; coordinator đã là role có trong kiến trúc sản phẩm. Cần xác định tutor
  là role mới hay tích hợp từ lecturer, và thêm vai trò phụ huynh/người giám hộ.
  \item \textbf{Cần thiết kế thêm:} hồ sơ nhu cầu học tập/điều chỉnh, tutor
  onboarding và xác minh, consent, yêu cầu ghép nối, thuật toán xếp hạng, màn
  hình điều phối, audit, lịch/feedback và quy trình báo sự cố.
  \item \textbf{Chưa thể tuyên bố đã có:} AI matching production, dữ liệu đào
  tạo, kết nối phụ huynh--học sinh, kiểm tra tutor chuyên biệt cho tự kỷ hoặc
  kết quả cải thiện học tập.
\end{itemize}

Vì vậy, đề xuất dưới đây là \textbf{kiến trúc mục tiêu để nghiên cứu và triển
khai sau này}, không phải sơ đồ mô tả tính năng đang chạy.

\textbf{AI-native trong phạm vi báo cáo này} nghĩa là app thiết kế ngay từ đầu
để biến yêu cầu học tập, tutor đã xác minh, lịch và phản hồi sau buổi học thành
dữ liệu có cấu trúc phục vụ matching; AI có thể giảm thời gian tìm shortlist
và giúp coordinator xem lý do. Nó không có nghĩa mọi tính năng phải gọi mô hình
hay AI được trao quyền tự ghép. Hiện chưa có bộ dữ liệu/outcome hoặc endpoint
matching trong source; vì vậy pilot bắt đầu bằng concierge và luật minh bạch,
sau đó mới đánh giá AI ở chế độ shadow.

\section{Module mục tiêu và luồng ghép nối}

\subsection{Sơ đồ module}

\begin{figure}[ht]
\centering
\resizebox{0.72\textwidth}{!}{%
\begin{tikzpicture}[
  font=\small,
  >=Latex,
  actor/.style={draw=ink,rounded corners=2mm,fill=white,minimum width=4.0cm,
    minimum height=1.05cm,align=center,text width=3.65cm,line width=0.8pt},
  client/.style={draw=teal,rounded corners=2mm,fill=sky,minimum width=8.5cm,
    minimum height=1.0cm,align=center,text width=8.0cm,line width=1pt},
  service/.style={draw=teal,rounded corners=2mm,fill=mint,minimum width=8.5cm,
    minimum height=0.95cm,align=center,text width=8.0cm,line width=1pt},
  human/.style={draw=ink,rounded corners=2mm,fill=amber,minimum width=8.5cm,
    minimum height=0.95cm,align=center,text width=8.0cm,line width=1pt},
  data/.style={draw=muted,rounded corners=2mm,fill=white,minimum width=8.5cm,
    minimum height=0.9cm,align=center,text width=8.0cm,dashed},
  flow/.style={-{Latex[length=2.5mm]},draw=teal,line width=0.9pt},
  feedback/.style={-{Latex[length=2mm]},draw=muted,dashed,line width=0.8pt}
]
\node[actor] (family) at (-2.7,6.2) {\textbf{Gia đình / học sinh}\\Mục tiêu học tập, cách học, lịch và consent};
\node[actor] (tutor) at (2.7,6.2) {\textbf{Tutor}\\Môn/khối lớp, kinh nghiệm, lịch và xác minh};
\node[client] (app) at (0,4.55) {\textbf{Ứng dụng Flutter / Web}\\Đăng ký nhu cầu, hồ sơ tutor, trạng thái và xác nhận};
\node[client] (api) at (0,2.95) {\textbf{Backend Matching Orchestrator}\\RBAC, consent, kiểm tra dữ liệu và tạo hồ sơ ghép nối};
\node[service] (filter) at (0,1.35) {\textbf{1. Điều kiện bắt buộc}\\Môn, khối lớp, lịch, hình thức, khu vực và tutor đã xác minh};
\node[service] (rank) at (0,-0.2) {\textbf{2. AI hỗ trợ xếp hạng}\\Chỉ hồ sơ đủ điều kiện; giải thích tiêu chí phù hợp và điểm chưa chắc chắn};
\node[human] (coordinator) at (0,-1.75) {\textbf{3. Coordinator duyệt}\\Kiểm tra an toàn, xem giải thích, đổi/thêm ứng viên và ghi lý do};
\node[client] (confirm) at (0,-3.3) {\textbf{Gia đình/học sinh xác nhận}\\Xem giới thiệu tutor, chọn buổi thử và có quyền yêu cầu ghép lại};
\node[client] (session) at (0,-4.85) {\textbf{Buổi học và phản hồi}\\Lịch học, trạng thái buổi, phản hồi theo consent và hỗ trợ ghép lại};
\node[data] (governance) at (0,-6.55) {\textbf{Xuyên suốt:} consent \textbar\ hồ sơ tối thiểu \textbar\ phân quyền
\textbar\ audit log \textbar\ lưu/xóa theo mục đích \textbar\ theo dõi phản hồi và sự cố};

\draw[flow] (family.south) -- (app.north west);
\draw[flow] (tutor.south) -- (app.north east);
\draw[flow] (app) -- (api);
\draw[flow] (api) -- (filter);
\draw[flow] (filter) -- (rank);
\draw[flow] (rank) -- (coordinator);
\draw[flow] (coordinator) -- (confirm);
\draw[flow] (confirm) -- (session);
\draw[feedback] (session) -- (governance);
\end{tikzpicture}%
}
\caption{Kiến trúc đề xuất. AI chỉ xếp hạng hồ sơ đủ điều kiện; coordinator
duyệt đề xuất; gia đình/học sinh xác nhận trước khi bắt đầu.}
\label{fig:modules}
\end{figure}

\clearpage
\subsection{Kiến trúc logic và ranh giới thành phần}

Sơ đồ module ở trên mô tả ai làm gì; sơ đồ sau mô tả các thành phần phần mềm
đề xuất và ranh giới dữ liệu. Luồng xử lý đi qua backend để xác thực quyền,
kiểm tra consent và ghi dấu quyết định. App mobile không gọi trực tiếp mô hình
AI hoặc gửi hồ sơ trẻ tới nhà cung cấp mô hình. Các kho dữ liệu trong hình là
thiết kế mục tiêu; chúng chưa có trong codebase.

\begin{figure}[ht]
\centering
\resizebox{\textwidth}{!}{%
\begin{tikzpicture}[
  font=\footnotesize,
  >=Latex,
  comp/.style={draw=ink,rounded corners=2mm,fill=sky,minimum width=3.7cm,
    minimum height=1.05cm,text width=3.35cm,align=center,line width=0.8pt},
  svc/.style={draw=teal,rounded corners=2mm,fill=mint,minimum width=3.7cm,
    minimum height=1.05cm,text width=3.35cm,align=center,line width=0.8pt},
  human/.style={draw=ink,rounded corners=2mm,fill=amber,minimum width=3.7cm,
    minimum height=1.05cm,text width=3.35cm,align=center,line width=0.8pt},
  store/.style={draw=muted,rounded corners=2mm,fill=white,minimum width=3.7cm,
    minimum height=1.05cm,text width=3.35cm,align=center,dashed,line width=0.8pt},
  flow/.style={-{Latex[length=2mm]},draw=teal,line width=0.8pt},
  logflow/.style={-{Latex[length=1.8mm]},draw=muted,dashed,line width=0.7pt}
]
\node[comp] (familyapp) at (-5.2,5.4) {\textbf{Flutter: gia đình / học sinh}\\Yêu cầu học, consent, lựa chọn tutor};
\node[comp] (tutorapp) at (-1.2,5.4) {\textbf{Flutter: tutor}\\Hồ sơ môn/lớp, lịch, nhận/từ chối ca};
\node[comp] (coordui) at (2.8,5.4) {\textbf{Web: coordinator}\\Shortlist, kiểm tra, duyệt/override};

\node[svc] (api) at (-1.2,3.45) {\textbf{Backend API + xác thực}\\Bearer/RBAC, kiểm tra chủ thể và quyền};
\node[svc] (consent) at (2.8,3.45) {\textbf{Consent + tối thiểu hóa}\\Chỉ mở yêu cầu khi đủ mục đích/quyền};
\node[svc] (matching) at (-1.2,1.45) {\textbf{Matching domain module}\\Tạo phiên bản yêu cầu và điều phối pipeline};
\node[store] (tutordata) at (2.8,1.45) {\textbf{Tutor registry}\\Xác minh, môn/lớp, lịch và phạm vi nhận ca};

\node[svc] (filter) at (-5.2,-0.55) {\textbf{Điều kiện bắt buộc}\\Môn/lớp, lịch, mode, địa bàn, xác minh};
\node[svc] (ranker) at (-1.2,-0.55) {\textbf{Ranker / AI worker}\\MVP: luật giải thích được; AI: nghiên cứu sau};
\node[svc] (reason) at (2.8,-0.55) {\textbf{Giải thích + thiếu dữ liệu}\\Tiêu chí, nguồn trường dữ liệu, độ bất định};
\node[human] (review) at (6.8,-0.55) {\textbf{Coordinator review queue}\\Duyệt, hỏi thêm, xếp lại hoặc không ghép};

\node[human] (choice) at (-1.2,-2.55) {\textbf{Gia đình/học sinh chọn}\\Đồng ý buổi thử hoặc đề nghị shortlist khác};
\node[comp] (booking) at (-1.2,-4.5) {\textbf{Booking / theo dõi}\\Buổi học, feedback, đổi tutor, báo sự cố};
\node[store] (profiles) at (2.8,-2.55) {\textbf{Kho dữ liệu có cấu trúc}\\Yêu cầu, hồ sơ tối thiểu, trạng thái ghép};
\node[store] (audit) at (6.8,-2.55) {\textbf{Consent + audit store}\\Ai xem/duyệt, lý do, phiên bản luật/model};

\draw[flow] (familyapp.south) -- (api.north west);
\draw[flow] (tutorapp.south) -- (api.north);
\draw[flow] (coordui.south) -- (api.north east);
\draw[flow] (api) -- (consent);
\draw[flow] (consent) -- (matching);
\draw[flow] (tutordata) -- (matching);
\draw[flow] (matching) -- (filter);
\draw[flow] (filter) -- (ranker);
\draw[flow] (ranker) -- (reason);
\draw[flow] (reason) -- (review);
\draw[flow] (review.south west) -- (choice.north east);
\draw[flow] (choice) -- (booking);
\draw[logflow] (matching.south east) -- (profiles.north);
\draw[logflow] (review.south) -- (audit.north);
\draw[logflow] (booking.east) -- (profiles.south west);
\draw[logflow] (booking.east) .. controls (3.4,-4.6) and (6.8,-4.6) .. (audit.south);
\end{tikzpicture}%
}
\caption{Kiến trúc logic mục tiêu. API là ranh giới xác thực; lọc loại hồ sơ
không đạt điều kiện trước khi ranker chạy; shortlist đi qua coordinator rồi
mới đến gia đình/học sinh. Kho dữ liệu và worker trong hình chưa được triển khai.}
\label{fig:target-architecture}
\end{figure}

\clearpage
\subsubsection{Các module và hợp đồng dữ liệu}

\begin{tabularx}{\textwidth}{p{2.8cm}p{5.1cm}Y}
\toprule
\textbf{Module} & \textbf{Input/output mục tiêu} & \textbf{Quyền và kiểm soát} \\
\midrule
Learner request & Môn, khối lớp, mục tiêu, lịch, địa bàn/mode, ngân sách và
điều chỉnh học tập do gia đình/học sinh chọn. & Guardian chỉ khai báo cho học
sinh mình quản lý; trường nhạy cảm là tùy chọn và không mặc định bắt buộc. \\
Tutor registry & Danh tính, môn/lớp, lịch trống, phạm vi địa bàn, kinh nghiệm
có thể xác minh, ngôn ngữ và điều chỉnh tutor tự nguyện cung cấp. & Coordinator
xác minh, hết hạn/xem lại hồ sơ; tutor tự cập nhật lịch và từ chối ca. \\
Eligibility filter & Tập tutor đã xác minh và còn nhận ca; trả về tập đủ điều
kiện hoặc lý do không có ứng viên. & Điều kiện bắt buộc do vận hành duyệt; không
để AI bỏ qua yêu cầu môn, lịch, an toàn hoặc quyền. \\
Ranker + explanation & Ứng viên hợp lệ, tiêu chí liên quan, điểm/tín hiệu thiếu
dữ liệu; trả shortlist có lý do truy ngược tới trường dữ liệu. & Không đưa chẩn
đoán, mức độ tự kỷ, ảnh, giọng nói hay suy luận cảm xúc vào ranking. \\
Coordinator review & Shortlist, lý do xếp hạng, ngoại lệ và khoảng trống thông
tin; trả quyết định duyệt/override/không ghép cùng lý do. & Coordinator giữ
quyền quyết định; mọi thay đổi trạng thái phải lưu người thao tác/thời điểm. \\
Consent and audit & Chủ thể, mục đích, phạm vi chia sẻ, thời hạn, lần rút
consent; sự kiện xem/sửa/đề xuất/duyệt/ghép lại. & Quyền theo vai trò, audit
không sửa, thời hạn lưu/xóa theo chính sách và rà soát pháp lý. \\
\bottomrule
\end{tabularx}

\clearpage
\subsubsection{Hợp đồng API tham khảo}

\begin{tabularx}{\textwidth}{p{5.0cm}p{3.2cm}Y}
\toprule
\textbf{Endpoint đề xuất (chưa tồn tại)} & \textbf{Ai gọi} &
\textbf{Kết quả và kiểm soát} \\
\midrule
\shortstack[l]{\textbf{POST}\\\texttt{/api/matching/requests}} &
Phụ huynh/người giám hộ hoặc coordinator
được phân quyền & Tạo yêu cầu tối thiểu; kiểm tra quyền quản lý học sinh và
consent trước khi lưu. \\
\shortstack[l]{\textbf{GET}\\\texttt{/api/matching/}\\\texttt{:id/shortlist}} &
Coordinator được giao
case & Trả candidates hợp lệ, tiêu chí, thiếu dữ liệu, phiên bản ranking;
không trả dữ liệu trẻ không cần thiết. \\
\shortstack[l]{\textbf{POST}\\\texttt{/api/matching/}\\\texttt{:id/review}} &
Coordinator & Lưu approve,
override, yêu cầu bổ sung hoặc no-match cùng lý do; không có thao tác “AI tự
ghép” bỏ qua người duyệt. \\
\shortstack[l]{\textbf{POST}\\\texttt{/api/matching/}\\\texttt{:id/choice}} &
Guardian / học sinh theo
quyền & Ghi đồng ý buổi giới thiệu, từ chối hoặc yêu cầu phương án khác; tutor
chỉ nhận thông tin tối thiểu đã được phép chia sẻ. \\
\shortstack[l]{\textbf{POST}\\\texttt{/api/sessions/}\\\texttt{:id/feedback}} &
Guardian, học sinh phù hợp độ
tuổi, tutor & Ghi feedback theo consent, yêu cầu đổi/dừng hoặc báo sự cố vào
luồng được phân quyền riêng. \\
\bottomrule
\end{tabularx}

\textbf{Triển khai theo nấc:} trong pilot concierge, coordinator có thể dùng
quy trình và hồ sơ được bảo vệ thay vì tạo model. MVP app gọi matching domain
trong backend hiện có; kho JSON hiện tại chỉ là persistence prototype, không
nên xem là nơi lưu mặc định cho dữ liệu nhạy cảm của trẻ. Trước khi có pilot
thật, cần chọn kho dữ liệu bền vững có kiểm soát truy cập, backup/xóa, mã hóa
và audit đã được rà soát. Chỉ tách worker/queue khi cần xử lý bất đồng bộ hoặc
quản trị model; không thêm cơ sở dữ liệu hay provider chỉ để gắn nhãn “AI”.

\subsection{Luồng vận hành đề xuất}

\begin{enumerate}
  \item \textbf{Tiếp nhận nhu cầu:} phụ huynh và học sinh chọn môn, khối lớp,
  mục tiêu, thời gian, online/trực tiếp, ngân sách và các điều chỉnh hữu ích
  cho buổi học. Cho phép bỏ qua thông tin chẩn đoán nếu không cần.
  \item \textbf{Kiểm tra consent và quyền truy cập:} người giám hộ cung cấp
  consent theo luồng được tư vấn pháp lý duyệt; học sinh được giải thích ở
  ngôn ngữ/cách thức phù hợp và có thể nêu ý kiến.
  \item \textbf{Xác minh tutor:} coordinator kiểm tra danh tính, năng lực theo
  môn/khối lớp, kinh nghiệm liên quan, tham chiếu và tiêu chuẩn bảo vệ trẻ em
  phù hợp quy trình vận hành.
  \item \textbf{Lọc điều kiện bắt buộc:} loại các hồ sơ sai môn, không phù hợp
  lịch/hình thức/địa bàn, chưa xác minh hoặc không thể nhận ca. AI không được
  dùng để bỏ qua điều kiện này.
  \item \textbf{Xếp hạng có giải thích:} hệ thống hiển thị vài hồ sơ và nêu
  yếu tố phù hợp như kinh nghiệm dạy đúng môn, lịch trùng, khả năng đáp ứng
  cách giao tiếp đã chọn. Mỗi yếu tố cần truy về dữ liệu hồ sơ cụ thể; không
  đưa ra kết luận lâm sàng hay suy đoán năng lực của trẻ.
  \item \textbf{Coordinator duyệt:} người điều phối xem cả ứng viên được đề
  xuất và lý do, có thể bỏ qua/đổi thứ tự, liên hệ hỏi thêm hoặc từ chối ghép.
  Hệ thống lưu người duyệt, thời điểm, phiên bản mô hình/luật và lý do override.
  \item \textbf{Gia đình/học sinh quyết định:} chỉ sau khi coordinator duyệt,
  gia đình nhận thông tin tối thiểu cần thiết để chọn buổi giới thiệu/thử.
  Không tự động đặt lịch hoặc chia sẻ toàn bộ hồ sơ cho tutor.
  \item \textbf{Theo dõi và ghép lại:} sau buổi đầu và theo mốc đã đồng ý,
  hỏi gia đình, học sinh và tutor về mức phù hợp; cho phép dừng hoặc yêu cầu đổi
  tutor mà không bị phạt vì phản hồi an toàn.
\end{enumerate}

\subsection{AI nên làm gì và không làm gì}

\begin{tabularx}{\textwidth}{p{3.2cm}Y Y}
\toprule
\textbf{AI có thể hỗ trợ} & \textbf{Quy tắc thiết kế} & \textbf{Không giao cho
AI} \\
\midrule
Tạo shortlist tutor đủ điều kiện; xếp hạng theo môn, lịch, modality, kinh
nghiệm dạy được xác minh và cách hỗ trợ học tập gia đình đã chọn. &
Ưu tiên dữ liệu có cấu trúc; nêu tiêu chí dẫn đến đề xuất; cho thấy thiếu dữ
liệu/độ bất định; coordinator có quyền override và ghi lý do. &
Chẩn đoán tự kỷ; đánh giá mức độ tự kỷ; quyết định học sinh có được nhận dịch
vụ hay không; quyết định tutor nào được nhận một em mà không qua người duyệt. \\
Tóm tắt yêu cầu dài thành checklist cho coordinator. &
Nếu dùng mô hình ngôn ngữ, chỉ tạo bản nháp để con người xác nhận; không gửi
hồ sơ sức khỏe, video hay ghi âm trẻ tới nhà cung cấp ngoài khi chưa có căn cứ,
consent và đánh giá dữ liệu. &
Suy diễn khả năng, hành vi hoặc mức độ rủi ro từ ảnh, giọng nói, cảm xúc hay
nhãn chẩn đoán; tự gửi tin nhắn/đề xuất cho trẻ. \\
\bottomrule
\end{tabularx}

\subsubsection{Đặc tả nghiên cứu ranking}

Chỉ xếp hạng sau khi áp dụng điều kiện bắt buộc. Một baseline minh bạch có thể
được biểu diễn:
\[
S(t,u)=w_sS_s+w_aS_a+w_mS_m+w_pS_p+w_eS_e,
\qquad \sum_i w_i=1.
\]
Trong đó $u$ là yêu cầu học tập, $t$ là hồ sơ tutor; các hạng lần lượt là
môn/lớp, lịch, hình thức/địa bàn, điều chỉnh đã chọn và kinh nghiệm đã xác
minh. Trọng số chưa được đo hay chốt; cần cùng gia đình, tutor và coordinator
xác định tiêu chí/ưu tiên,
ghi rõ trade-off và kiểm tra xem dữ liệu có chứng minh được tiêu chí hay không.
Không đưa tình trạng tự kỷ hoặc chẩn đoán làm “điểm phù hợp”; ưu tiên nhu cầu
thực tế và sự điều chỉnh do người học chọn. Nếu thiếu dữ liệu quan trọng, trả
“chưa đủ căn cứ” hoặc no-match để coordinator hỏi thêm, thay vì đoán.

\begin{tabularx}{\textwidth}{p{3.1cm}p{5.0cm}Y}
\toprule
\textbf{Chỉ số pilot} & \textbf{Định nghĩa đo} & \textbf{Cách diễn giải} \\
\midrule
Thời gian tới shortlist duyệt & Trung vị và p90 từ lúc nhận yêu cầu đủ consent
đến khi coordinator có shortlist chấp nhận được. & Đo hiệu suất vận hành; không
phải chất lượng học tập. Báo riêng case chờ tutor/thiếu dữ liệu. \\
Tỷ lệ có tutor nhận ca & Yêu cầu có ít nhất một tutor đủ điều kiện, được đề
xuất và xác nhận nhận buổi đầu / tổng yêu cầu đủ điều kiện. & Phân rã theo
môn, lịch, địa bàn, hình thức và lý do không ghép; đừng che no-match bằng cách
loại ca khó khỏi mẫu. \\
Tỷ lệ override & Shortlist coordinator đổi thứ tự, bỏ ứng viên hoặc yêu cầu
bổ sung / tổng shortlist được xem. & Tỷ lệ cao không mặc định là xấu; phân loại
lý do để tìm thiếu dữ liệu hoặc tiêu chí chưa đúng. \\
Chấp nhận và buổi đầu & Gia đình chọn ít nhất một ứng viên; sau đó theo dõi
buổi đầu được đặt/hoàn tất, đổi tutor, dừng dịch vụ. & Chấp nhận hồ sơ khác với
chất lượng tutor hay hiệu quả học tập; báo các outcome riêng. \\
Công bằng và an toàn & So sánh no-match, thời gian chờ, override, chấp nhận và
sự cố theo nhóm đủ cỡ mẫu và có consent phù hợp. & Không kết luận parity từ
nhóm quá nhỏ; không công bố phân nhóm có thể tái nhận diện trẻ. \\
\bottomrule
\end{tabularx}

MVP nên bắt đầu bằng lọc điều kiện và bảng điểm luật minh bạch, có thể giải
thích từng tiêu chí. Giai đoạn sau mới so sánh mô hình xếp hạng với quyết định
của coordinator ở chế độ shadow. Chỉ cân nhắc học từ dữ liệu outcome khi có
đủ consent, chất lượng nhãn, kiểm tra sai lệch theo nhóm và đánh giá tác động
đến quyền trẻ em. Không có ngưỡng số mẫu được khẳng định trong báo cáo này;
ngưỡng phải được xác định cùng chuyên gia dữ liệu và chuyên gia giáo dục.

\section{Quyền riêng tư, an toàn và quản trị}

Thông tin về chẩn đoán, điều chỉnh học tập, hành vi hoặc sức khỏe có thể nhạy
cảm và gắn với trẻ em. Thiết kế cần mặc định thu thập ít nhất có thể: thường
ưu tiên mục tiêu học tập, cách giao tiếp, lịch, sở thích và điều chỉnh mà gia
đình tự chọn, thay vì lưu hồ sơ y tế đầy đủ. Cần xác định ai được xem từng
trường dữ liệu, mục đích sử dụng, thời hạn giữ, cách sửa/xóa, quy trình khi rút
consent và ứng phó khi có sự cố.

Tại ngày lập báo cáo, Luật Bảo vệ dữ liệu cá nhân số 91/2025/QH15 và Nghị định
356/2025/NĐ-CP đã có hiệu lực từ 01/01/2026; Luật Trí tuệ nhân tạo số
134/2025/QH15 có hiệu lực từ 01/03/2026 \cite{pdpl2025,decree356,ailaw2025}.
Trước pilot, cần luật sư/chuyên gia tuân thủ xác nhận vai trò xử lý dữ liệu,
consent với trẻ và người đại diện, thông báo, chia sẻ với nhà cung cấp AI,
đánh giá tác động và nghĩa vụ cụ thể theo văn bản hiện hành. Báo cáo này không
thay thế ý kiến pháp lý.

UNICEF khuyến nghị hệ thống AI hướng đến trẻ em phải quan tâm an toàn, riêng
tư, công bằng, khả năng giải thích, trách nhiệm giải trình, quyền trẻ em và
giám sát \cite{unicef2025ai}. Các điểm chuyển thành yêu cầu sản phẩm:

\begin{itemize}
  \item Không dùng thông tin tự kỷ như một biến điểm làm giảm cơ hội được ghép.
  Hãy mô tả nhu cầu hỗ trợ cụ thể và cho gia đình quyền chỉnh/sửa.
  \item Tách thông tin nhận dạng khỏi dữ liệu dùng để đánh giá matching; không
  dùng nội dung hồ sơ để train model ngoài mục đích đã thông báo/cho phép.
  \item Không để AI là người liên hệ trực tiếp với trẻ. Giao tiếp về matching
  đi qua người giám hộ và coordinator theo quy trình.
  \item Có nút báo lo ngại, dừng ghép, liên hệ coordinator, sửa hồ sơ và yêu
  cầu xem lại. Sự cố bảo vệ trẻ em phải có luồng xử lý riêng.
  \item Đánh giá chênh lệch về số đề xuất, tỷ lệ được nhận và tỷ lệ ghép thành
  công; không lấy một chỉ số tổng thể làm bằng chứng công bằng.
\end{itemize}

\section{Kế hoạch kiểm chứng và tiêu chí quyết định}

Khung sau là \textbf{kế hoạch lấy dữ liệu đề xuất}, không phải số liệu đã thu:
phỏng vấn định tính 15--20 gia đình/người giám hộ, 10--15 tutor, và 3--5
đại diện trung tâm/trường/coordinator tại một địa bàn; mời học sinh tham gia
theo cách phù hợp độ tuổi, giao tiếp, consent và mong muốn của các em. Chọn mẫu
có chủ đích để khám phá rào cản và tiêu chí, không tuyên bố đại diện dân số.
Pilot concierge gợi ý 20--30 yêu cầu ghép trong 8--12 tuần; nếu mẫu nhỏ hoặc
chỉ đến từ một đối tác, kết quả chỉ có tính mô tả và không chứng minh quy mô
quốc gia hay hiệu quả học tập.

\begin{tabularx}{\textwidth}{p{2.55cm}p{3.15cm}Y}
\toprule
\textbf{Giai đoạn} & \textbf{Việc làm} & \textbf{Bằng chứng/đầu ra} \\
\midrule
Discovery &
Phỏng vấn định tính theo khung 15--20 gia đình, 10--15 tutor, 3--5 đối tác/
coordinator; ít nhất một địa bàn. &
Bản đồ cách tìm tutor hiện tại; 5--7 tiêu chí ghép thực tế; nhu cầu về riêng
tư/an toàn; ai trả tiền; phản ứng với thẻ giá giả định; danh sách nguồn tutor
có thể xác minh. Không tính là khảo sát đại diện. \\
\addlinespace
Concierge pilot &
Coordinator tự quản lý khoảng 20--30 yêu cầu ghép trong 8--12 tuần; không cần
AI ra quyết định. &
Tỷ lệ chấp nhận đề xuất; số ngày đến shortlist/buổi đầu; buổi hoàn tất/hủy;
tỷ lệ ghép lại; phút coordinator cho mỗi ca; chi phí và willingness-to-pay;
mọi sự cố/an toàn. \\
\addlinespace
Shadow ranking &
Chạy ranking trên bản sao dữ liệu tối thiểu, không hiển thị kết quả cho gia
đình và không tác động quyết định. &
So khớp với shortlist coordinator; lý do chênh lệch; thiếu dữ liệu; kiểm tra
phân phối kết quả theo nhóm; đánh giá khả năng giải thích và an toàn. \\
\addlinespace
Thử nghiệm có kiểm soát &
Chỉ khi đã qua privacy/child-rights review; coordinator duyệt mọi shortlist;
gia đình/học sinh vẫn chọn. &
So sánh với concierge baseline về thời gian và mức phù hợp, kiểm tra khiếu
nại, đổi tutor và các khác biệt giữa nhóm; dừng nếu xuất hiện rủi ro không
kiểm soát được. \\
\bottomrule
\end{tabularx}

\subsection{Câu hỏi go/no-go}

\begin{enumerate}
  \item Gia đình có vấn đề tìm tutor đủ lớn để thử dịch vụ, và có đồng ý cho
  một coordinator trung gian hỗ trợ không?
  \item Có tutor đạt chuẩn và đủ lịch cho \emph{từng môn/địa bàn} hay chỉ có
  nhiều hồ sơ đăng ký nhưng không thể nhận ca?
  \item Coordinator có thể phục vụ số ca cần thiết với chất lượng và chi phí
  bền vững không?
  \item Gia đình, tutor hoặc đối tác có sẵn lòng trả tiền cho phần điều phối,
  xác minh và hỗ trợ đổi tutor không?
  \item Matching dựa trên mục tiêu/cách hỗ trợ mà gia đình chọn có được chấp
  nhận hơn ghép theo môn và lịch đơn thuần không?
  \item AI có cải thiện thời gian tìm shortlist mà không làm giảm quyền lựa
  chọn, niềm tin, an toàn hoặc công bằng không?
\end{enumerate}

\textbf{Điều kiện mở rộng:} cần thấy cả nhu cầu lặp lại, nguồn cung tutor đã
xác minh, kinh tế đơn vị không âm sau khi tính đủ lao động điều phối, cùng quy
trình privacy/safeguarding có thể kiểm tra. Chưa nên scale chỉ vì tỷ lệ chấp
nhận top-k cao hoặc vì mô hình AI chạy được.

\section{Rủi ro chính}

\begin{tabularx}{\textwidth}{p{3.1cm}Y Y}
\toprule
\textbf{Rủi ro} & \textbf{Hậu quả có thể xảy ra} & \textbf{Cách giảm thiểu} \\
\midrule
Thiếu mật độ cung-cầu &
Không có đủ tutor đúng môn/lịch tại từng địa bàn; trải nghiệm chờ lâu. &
Bắt đầu ở một khu vực/đối tác có nguồn cung thật; nhận đăng ký có kiểm soát;
đưa lựa chọn online/chờ danh sách với consent rõ. \\
Thiên lệch từ dữ liệu thiếu &
Ranking ưu tiên hồ sơ phổ biến, ngôn ngữ/địa bàn hoặc nhóm đã có nhiều ca;
đề xuất không công bằng cho nhu cầu ít gặp. &
Giữ điều kiện minh bạch; hiển thị thiếu dữ liệu; theo dõi kết quả phân nhóm;
cho coordinator override; không train sớm trên dataset nhỏ. \\
Kỳ vọng quá mức về AI &
Gia đình hiểu ranking là bảo chứng tutor phù hợp hoặc kết quả giáo dục. &
Gọi đúng là gợi ý; giải thích giới hạn; coordinator trao đổi kỳ vọng và kiểm
tra sau buổi đầu. \\
Rò rỉ dữ liệu trẻ em &
Thông tin nhạy cảm bị xem/chia sẻ/lưu lâu hơn mục đích. &
Tối thiểu hóa dữ liệu, quyền theo vai trò, audit, thời hạn xóa, consent phù
hợp và thẩm định luồng nhà cung cấp trước khi bật AI. \\
Kinh tế đơn vị không đủ &
Chi phí coordinator vượt phí hoa hồng; tutor/gia đình giao dịch ngoài nền
tảng sau khi được giới thiệu. &
Đo phút/ca và tỷ lệ duy trì; định giá phần điều phối có giá trị; thử hợp đồng
với tổ chức trước khi mở rộng B2C. \\
\bottomrule
\end{tabularx}

\section{Kết luận}

Ý tưởng có cơ sở để \textbf{thử như một dịch vụ kết nối tutor do con người điều
phối}, tập trung vào đúng môn, lịch, sở thích học tập và hỗ trợ cá nhân mà gia
đình/học sinh lựa chọn. Các số liệu về tự kỷ và khuyết tật tạo bối cảnh cho
nhu cầu hòa nhập, nhưng không xác định số khách hàng trả phí. Vì thế, viability
thương mại vẫn là câu hỏi cần trả lời bằng discovery và pilot tại địa bàn cụ
thể.

Vai trò coordinator là một phần thiết yếu của dịch vụ: xác minh tutor, hiểu
ngữ cảnh, xem lại shortlist AI, xin xác nhận của gia đình/học sinh, theo dõi
buổi học và xử lý đổi ghép. AI chỉ nên xếp hạng các hồ sơ đã qua điều kiện bắt
buộc, giải thích vì sao, báo thiếu dữ liệu và chịu kiểm soát bởi con người.
Đây là ranh giới vừa hợp với sản phẩm trẻ em vừa làm rõ giá trị của app: không
chỉ tìm người theo môn, mà tổ chức một quy trình ghép nối có thể giải thích,
điều chỉnh và chịu trách nhiệm.

\section*{Ghi chú phương pháp}
\addcontentsline{toc}{section}{Ghi chú phương pháp}

Đây là desk research được mở rộng trong vòng này, chốt ngày 02/10/2026. Số liệu được giữ nguyên
phạm vi tuổi/địa bàn/method do nguồn công bố. Không có phỏng vấn khách hàng,
khảo sát willingness-to-pay, dữ liệu vận hành hay chi phí kế toán được cung
cấp trong repository; mọi phần kinh doanh phụ thuộc những dữ liệu này đều là
giả thuyết hoặc kịch bản minh họa. Các chức năng AI trong báo cáo là thiết kế
đề xuất, chưa được xác nhận đã triển khai.

\begin{thebibliography}{99}
\footnotesize
\setlength{\itemsep}{1pt}
\setlength{\parskip}{0pt}
\setlength{\parsep}{0pt}

\bibitem{vui2022}
Le Thi Vui, Duc Duong Minh, Nguyen Thuy Quynh, Nguyen Thi Huong Giang,
Vu Thi Thanh Mai, Bui Thi Thu Ha, \& Hoang Van Minh.
(2022).
\newblock Early screening and diagnosis of autism spectrum disorders in Vietnam:
A population-based cross-sectional survey.
\newblock \emph{Journal of Public Health Research, 11}(2), 2460.
\newblock \href{https://doi.org/10.4081/jphr.2021.2460}{doi:10.4081/jphr.2021.2460}.

\bibitem{tran2024}
Tran, T.~T., Nguyen, V.~T., Nguyen, M.~P., Huynh Nguyen, P.~Q., Le, T.~T.~H.,
Nguyen, H.~Y., Nguyen, T.~H., Nguyen, T.~L., Le, C.~T., \& Nguyen, V.~T.
(2024).
\newblock Prevalence of autism spectrum disorder diagnosed according to DSM-5
criteria and associated factors in preschoolers in southern Vietnam.
\newblock \emph{La Clinica Terapeutica, 175}(6), 405--411.
\newblock \href{https://doi.org/10.7417/CT.2024.5147}{doi:10.7417/CT.2024.5147}.

\bibitem{who2025}
World Health Organization.
(2025, September 17).
\newblock Autism.
\newblock \url{https://www.who.int/westernpacific/newsroom/fact-sheets/detail/autism-spectrum-disorders}.

\bibitem{unicef2025inclusive}
United Nations Children's Fund.
(2025, June).
\newblock Viet Nam: Making inclusive education a reality for every child.
\newblock In \emph{Core Resources Impact Companion: For Every Child, Everywhere,
Annual Report 2024}.
\newblock \url{https://www.unicef.org/media/171656/file/UNICEF-Core-Resources-Annual-Report-2024-Impact-Companion.pdf}.

\bibitem{unicef2025ai}
United Nations Children's Fund, Office of Strategy and Evidence -- Innocenti.
(2025, December).
\newblock \emph{Guidance on AI and children: Version 3.0}.
\newblock \url{https://www.unicef.org/innocenti/reports/policy-guidance-ai-children}.

\bibitem{pdpl2025}
Quốc hội nước Cộng hòa xã hội chủ nghĩa Việt Nam.
(2025).
\newblock Luật Bảo vệ dữ liệu cá nhân số 91/2025/QH15, hiệu lực 01/01/2026.
\newblock \url{https://chinhphu.vn/?docid=214590&pageid=27160&typegroupid=3}.

\bibitem{decree356}
Chính phủ nước Cộng hòa xã hội chủ nghĩa Việt Nam.
(2025).
\newblock Nghị định 356/2025/NĐ-CP hướng dẫn Luật Bảo vệ dữ liệu cá nhân,
hiệu lực 01/01/2026.
\newblock \url{https://vanban.chinhphu.vn/?classid=1&docid=216387&pageid=27160}.

\bibitem{ailaw2025}
Quốc hội nước Cộng hòa xã hội chủ nghĩa Việt Nam.
(2025).
\newblock Luật Trí tuệ nhân tạo số 134/2025/QH15, hiệu lực 01/03/2026.
\newblock \url{https://vanban.chinhphu.vn/?classid=1&docid=216334&pageid=27160}.

\bibitem{cdc2025}
Kelly A. Shaw et al.
(2025).
\newblock Prevalence and Early Identification of Autism Spectrum Disorder Among
Children Aged 4 and 8 Years --- Autism and Developmental Disabilities Monitoring
Network, 16 Sites, United States, 2022.
\newblock \emph{Morbidity and Mortality Weekly Report, 74}(2), 1--22.
\newblock \url{https://www.cdc.gov/mmwr/volumes/74/ss/ss7402a1.htm}.

\bibitem{gso2023}
Tổng cục Thống kê Việt Nam.
(2024).
\newblock Thông cáo báo chí về kết quả Điều tra người khuyết tật năm 2023.
\newblock \url{https://www.nso.gov.vn/tin-tuc-thong-ke/2024/11/thong-cao-bao-chi-ve-ket-qua-dieu-tra-nguoi-khuyet-tat-nam-2023/}.

\bibitem{unescoai2025}
UNESCO.
(2025, October 24).
\newblock \emph{Viet Nam: Artificial Intelligence Readiness Assessment Report}.
\newblock \url{https://articles.unesco.org/sites/default/files/medias/fichiers/2025/10/Viet%20Nam%20Artificial%20Intelligence%20Readiness%20Assessment%20Report.pdf}.

\bibitem{unescoedtech2026}
UNESCO SDG 4 Knowledge Hub.
(2026).
\newblock Advancing equitable and inclusive digital transformation in education
in Viet Nam (program period 2022--2026).
\newblock \url{https://www.unesco.org/sdg4education2030/en/knowledge-hub/advancing-equitable-and-inclusive-digital-transformation-education-viet-nam?hub=25}.

\bibitem{worldbankdeaf2023}
Noah Yarrow, Lauri Pynnonen, Chuyu Song, Riaz Bhardwaj, \& Mario Spiezio.
(2023).
\newblock \emph{Use of Assistive Education Technologies to Support Children
with Visual and Hearing Difficulties in the East Asia and Pacific Region}.
\newblock World Bank.
\newblock \url{https://documents1.worldbank.org/curated/en/099062923221035313/pdf/P17621707c0e430980bc600f762c5761c85.pdf}.

\bibitem{spmarchitecture2026}
SPM Mobile project.
(2026, October 2).
\newblock \emph{Kiến trúc hệ thống SPM}.
\newblock Tài liệu nội bộ trong repository, \texttt{docs/architecture/system-architecture.md},
đối chiếu mã nguồn ngày 02/10/2026.

\end{thebibliography}

\end{document}
